The isomorphism $B_\gamma \cong \Mat_2(\varepsilon_+ B_\gamma \varepsilon_+) = \Mat_2(\varepsilon_+ A) \cong \Mat_2(\Cl^0(p,q))$ chains Lemmas \claimref{transport.matrix-units-peirce-isomorphism} and \claimref{transport.corner-even-subalgebra}; note the matrix units and the corner are the same for both signs of $\gamma$ up to the harmless factor $\gamma$ in $E_{21}$, which is why the right-hand side is $\gamma$-independent; this proves (3). For (1): $\varepsilon_+ A$ is semisimple (a direct summand of the semisimple $A$, Theorem \claimref{transport.classification-real-clifford-algebras}) with centre $\varepsilon_+ Z(A) = \R\,\varepsilon_+$ one-dimensional (here $Z(A) = \R \oplus \R\omega = \R\varepsilon_+ \oplus \R\varepsilon_-$ by Lemma \claimref{transport.centre}, and the centre of a direct sum is the direct sum of the centres), and a semisimple algebra whose centre is one-dimensional has exactly one Wedderburn block, i.e.\ is simple. Matrix algebras over simple algebras are simple (\claimref{std-algebra.morita-amplification}), so $B_\gamma$ is simple; its centre is $\R$ directly by Lemma \claimref{transport.centre}, since $n + 1$ is even. (2) is Proposition \claimref{transport.ascent-renders-grade-involution-inner} together with the identification of $\varepsilon_\pm$ as $E_{11}, E_{22}$; a diagonal matrix unit in $\Mat_2$ of a unital algebra is never central.
